\documentclass[11pt]{article}

\usepackage[margin=1in]{geometry}
\usepackage[T1]{fontenc}
\usepackage[utf8]{inputenc}
\usepackage{amsmath,amssymb,amsfonts,mathtools}
\usepackage{booktabs}
\usepackage{array}
\usepackage{longtable}
\usepackage{enumitem}
\usepackage{hyperref}

\hypersetup{
  colorlinks=true,
  linkcolor=blue,
  urlcolor=blue,
  citecolor=blue
}

\setlist[itemize]{leftmargin=1.5em}
\setlist[enumerate]{leftmargin=1.8em}

\newcommand{\R}{\mathbb{R}}
\newcommand{\N}{\mathbb{N}}
\newcommand{\1}{\mathbf{1}}
\newcommand{\X}{\mathbf{X}}
\newcommand{\y}{\mathbf{y}}
\newcommand{\E}{\mathbb{E}}
\newcommand{\Var}{\mathrm{Var}}
\newcommand{\argmin}{\mathop{\mathrm{argmin}}}
\newcommand{\loss}{\mathcal{L}}
\newcommand{\D}{\mathcal{D}}
\newcommand{\F}{\mathcal{F}}
\newcommand{\T}{\mathcal{T}}
\newcommand{\seq}{\mathcal{S}}

\title{Axiomatic and Formal Description of the CatBoost + CNN Monthly Return Pipeline}
\author{ }
\date{ }

\begin{document}
\maketitle

\begin{abstract}
We study a supervised learning problem on a monthly company-level panel. The objective is to predict the monthly gross return of a company using a two-stage model: a tabular first-stage learner based on CatBoost and a temporal residual learner based on a one-dimensional convolutional neural network. The pipeline is leakage-safe, chronological, and split-aware. This document states the problem axiomatically, defines the preprocessing operators mathematically, maps the code to formal transformations, presents pseudocode, and then develops CatBoost and CNN from first principles in appendices.
\end{abstract}

\tableofcontents
\newpage

\section{Problem Statement}

We consider a panel of monthly company observations. Each observation belongs to exactly one company and exactly one month. The prediction target is the monthly gross return. The goal is to construct a function that maps the information available at month $t$ for company $j$ to an estimate of the realized return at that month.

Let $j \in \{1,\dots,J\}$ index companies and $t \in \{1,\dots,T_j\}$ index months available for company $j$. Let the realized return be
\[
y_{j,t} \in \R.
\]
Let the raw observation vector be denoted by $r_{j,t}$ and let the engineered feature vector be denoted by $x_{j,t}$. The learning objective is to approximate the conditional relationship
\[
y_{j,t} \approx f(x_{j,t}),
\]
where $f$ is not a single function in this pipeline, but a composition
\[
f(x_{j,t}) = f_{\mathrm{CB}}(x_{j,t}) + f_{\mathrm{CNN}}(\seq_{j,t}).
\]

The first term captures structured tabular and cross-sectional signal. The second term captures remaining temporal structure in the residuals.

\section{Axiomatic Framework}

We state the pipeline under a small axiom system.

\subsection{Axiom 1: Panel Identity}

Every row corresponds to one and only one pair $(j,t)$ after duplicate collapse. Formally, there exists an injective mapping from the processed table rows to company-month keys.

\subsection{Axiom 2: Chronological Order}

For each fixed company $j$, the months are totally ordered by time:
\[
t_1 < t_2 \implies \text{month}(t_1) \prec \text{month}(t_2).
\]
All lagged and rolling constructions must respect this order.

\subsection{Axiom 3: No Future Leakage}

Every feature at time $t$ must be a deterministic function of data available at or before $t$. If $\phi$ denotes the feature construction map, then
\[
x_{j,t} = \phi\big(\{r_{j,\tau} : \tau \le t\}\big)
\]
and never depends on $r_{j,\tau}$ for $\tau > t$.

\subsection{Axiom 4: Split Integrity}

The train, validation, and test splits are disjoint in time. If $\T_{\mathrm{train}}, \T_{\mathrm{val}}, \T_{\mathrm{test}}$ denote the month sets, then
\[
\T_{\mathrm{train}} \cap \T_{\mathrm{val}} = \varnothing,\quad
\T_{\mathrm{val}} \cap \T_{\mathrm{test}} = \varnothing.
\]
In the implemented pipeline, train ends before validation begins, and validation ends before test begins.

\subsection{Axiom 5: Deterministic Canonicalization}

Whenever duplicate or malformed observations are present, the preprocessing stage reduces them to a unique canonical representation using deterministic rules. This ensures repeatability and stable feature construction.

\section{Data Model and Notation}

We define the raw panel as
\[
\D^{\mathrm{raw}} = \{ (j,t,r_{j,t}, m_{j,t}, c_{j,t}) \},
\]
where:
\begin{itemize}
  \item $r_{j,t}$ denotes the target-return-bearing row data,
  \item $m_{j,t}$ denotes numeric and mixed-type metadata,
  \item $c_{j,t}$ denotes categorical labels.
\end{itemize}

The processed panel is denoted by
\[
\D = \{(j,t,x_{j,t},y_{j,t})\}.
\]
The final feature map is a deterministic transformation
\[
\phi : \D^{\mathrm{raw}} \to \R^d.
\]
The sequence map is
\[
\psi : \D \to \R^{d \times L},
\]
where $L$ is the sequence length used by the CNN.

\section{Approach Overview}

We use a two-stage residual learning strategy.

\subsection{Stage 1: CatBoost}

CatBoost learns a function $f_{\mathrm{CB}} : \R^d \to \R$ that predicts the target from tabular features. It captures nonlinear cross-sectional dependencies, missingness structure, and categorical interactions.

\subsection{Stage 2: CNN on Residuals}

The second model learns a correction term from a causal sequence of past features. If the first-stage prediction is $\hat y^{\mathrm{CB}}_{j,t}$, then the residual is
\[
e_{j,t} = y_{j,t} - \hat y^{\mathrm{CB}}_{j,t}.
\]
The CNN learns a mapping from a sequence tensor to this residual.

\subsection{Final Prediction}

The combined predictor is
\[
\hat y_{j,t} = \hat y^{\mathrm{CB}}_{j,t} + \hat e^{\mathrm{CNN}}_{j,t}.
\]
This decomposition is the central modeling hypothesis of the notebook.

\section{Code-to-Mathematics Correspondence}

This section formalizes the major code blocks in the notebook. The objective is not merely to narrate what the code does, but to state the exact mathematical operator induced by each block.

\subsection{Reading a Split}

The function \texttt{read\_split(path, split\_name)} reads a CSV file and appends a split label.

Mathematically, if $\D_s$ is the set of rows in split $s \in \{\mathrm{train},\mathrm{validation},\mathrm{test}\}$, then
\[
\D_s = \{(u, s) : u \in \text{CSV}(s)\}.
\]
The split label is an exogenous tag used for bookkeeping and later filtering.

\subsection{Dropping Previously Engineered Features}

The function \texttt{drop\_previous\_engineering} removes stale engineered columns if they already exist. Let $C$ be the full column set and let $C_{\mathrm{drop}}$ be the set of columns to remove. Then the operator is the projection
\[
\pi_{C\setminus C_{\mathrm{drop}}}(\D).
\]
This projection prevents feature duplication and prevents recursive application of the same engineering rule from stacking derived variables on top of derived variables.

\subsection{Type Coercion}

The function \texttt{coerce\_panel\_types} performs three canonical transformations.

\paragraph{Date coercion.}
The month field is mapped to a time element in a totally ordered time set. Let $\tau : \text{strings} \to \mathcal{T}$ be the conversion into timestamps.

\paragraph{Identifier coercion.}
The company identifier is mapped into a nullable integer space. This does not change the entity identity; it only standardizes representation.

\paragraph{Target coercion.}
The return column is mapped into $\R$.

For every column $q$ not in the exempted set $E = \{\texttt{co\_code}, \texttt{Month}, \texttt{monthly\_gross\_return}, \texttt{source\_split}\} \cup C_{\mathrm{cat}}$, we apply
\[
q \mapsto \nu(q),
\]
where $\nu$ is numeric coercion with missing-value propagation.

\subsection{Duplicate Collapse}

The function \texttt{collapse\_duplicate\_company\_months} enforces the uniqueness of the key $(j,t)$.

Define an equivalence relation on rows by
\[
u \sim v \iff (j_u,t_u) = (j_v,t_v).
\]
Each equivalence class corresponds to one company-month key. The code selects the last representative of each class. If $[u]$ is an equivalence class ordered by appearance, then the canonical representative is
\[
\rho([u]) = \max_{\prec} [u],
\]
where $\prec$ is the input ordering. The resulting processed panel satisfies key uniqueness.

\subsection{Lag Feature Construction}

The function \texttt{engineer\_features} constructs lagged returns. Let $y_{j,t}$ be the target return. For lag $\ell \in \{1,3,6\}$,
\[
\mathrm{lag\_ret}_{\ell}(j,t) = y_{j,t-\ell}.
\]
This is a shift operator on the time index.

If \texttt{log\_ret} exists, then for the same lags,
\[
\mathrm{log\_ret\_lag}_{\ell}(j,t) = \log\_\mathrm{ret}(j,t-\ell).
\]
These are causal by construction because they only use earlier months.

\subsection{Rolling Feature Construction}

Let $\bar y_{j,t}^{(w)}$ denote the rolling mean with window $w \in \{3,6\}$. The code first shifts the target by one month and then computes a rolling statistic on the shifted series. Hence
\[
\bar y_{j,t}^{(w)} = \frac{1}{|I_{j,t}^{(w)}|} \sum_{\tau \in I_{j,t}^{(w)}} y_{j,\tau},
\]
where $I_{j,t}^{(w)}$ is the set of the previous $w$ observations ending at $t-1$, truncated when fewer than $w$ historical values exist.

Similarly, the rolling standard deviation is
\[
s_{j,t}^{(w)} = \sqrt{\frac{1}{|I_{j,t}^{(w)}| - 1} \sum_{\tau \in I_{j,t}^{(w)}} (y_{j,\tau} - \bar y_{j,t}^{(w)})^2}
\]
when at least two observations exist; otherwise the sample variance semantics inherited from the implementation apply.

\subsection{Month Index}

The month index is a deterministic calendar embedding:
\[
\mathrm{month\_index}(j,t) = 12 \cdot \mathrm{year}(j,t) + \mathrm{month}(j,t).
\]
This is a monotone index on the calendar time axis.

\subsection{Missingness Indicators}

For each feature column $f$, the indicator is
\[
\mathrm{is\_missing}_f(j,t) =
\begin{cases}
1, & \text{if } f(j,t) \text{ is missing},\\
0, & \text{otherwise}.
\end{cases}
\]
This is a binary feature measuring the existence of information, not its magnitude.

\subsection{Feature Set Filter}

The feature set excludes identifiers, the target, split labels, and internal outputs. Let $C$ be the total column set and let $E$ be the excluded set. Then the model feature set is
\[
C_{\mathrm{feat}} = C \setminus E.
\]
The categorical subset is
\[
C_{\mathrm{cat}}^{\mathrm{feat}} = C_{\mathrm{feat}} \cap C_{\mathrm{cat}}.
\]

\subsection{Chronological Split Verification}

Let $\max \T_{\mathrm{train}}$ denote the latest train month, $\min \T_{\mathrm{val}}$ the earliest validation month, and $\min \T_{\mathrm{test}}$ the earliest test month. The code enforces
\[
\max \T_{\mathrm{train}} < \min \T_{\mathrm{val}} < \min \T_{\mathrm{test}}.
\]
This guarantees that the training process sees only past data relative to validation and test.

\subsection{CatBoost Feature Frame}

The function \texttt{make\_feature\_frame} reindexes rows to $C_{\mathrm{feat}}$ and normalizes categorical entries as strings. This is a canonicalization map
\[
\chi : \D \to \R^d \times \Sigma^{|C_{\mathrm{cat}}^{\mathrm{feat}}|},
\]
where $\Sigma$ is the categorical alphabet.

\subsection{CatBoost Pool and Training}

The function \texttt{make\_pool} constructs the supervised learning pair
\[
\big( X, y \big) = \big( \X_{\mathrm{train}}, \y_{\mathrm{train}} \big)
\]
with categorical feature indices specified explicitly.

Training solves an empirical risk minimization problem of the form
\[
\hat f_{\mathrm{CB}} = \argmin_{f \in \mathcal{F}_{\mathrm{CB}}} \sum_{(x_i,y_i) \in \X_{\mathrm{train}}} \ell(y_i, f(x_i)),
\]
with $\ell(y,\hat y) = (y-\hat y)^2$ in the implementation.

\subsection{Residual Construction}

For each non-training row, the base residual is
\[
r_{j,t}^{\mathrm{base}} = y_{j,t} - \hat y^{\mathrm{CB}}_{j,t}.
\]
The training split assigns no residual supervision to the CNN stage. Validation and test residuals are the OOS targets.

\subsection{CNN Sequence Split}

The validation split is further partitioned into a CNN-train subset and a CNN-holdout subset. Let $\T_{\mathrm{val}} = \T_{\mathrm{cnn\_train}} \cup \T_{\mathrm{cnn\_holdout}}$ with disjoint union. This creates a second-level chronological separation.

\subsection{Encoding for the CNN}

The function \texttt{encode\_for\_cnn} maps categorical columns to integer codes and numeric columns to floating-point values. Formally, for each column $f$:
\[
\eta_f(x) =
\begin{cases}
\mathrm{code}(x), & f \in C_{\mathrm{cat}}^{\mathrm{feat}},\\
\mathrm{to\_float}(x), & f \notin C_{\mathrm{cat}}^{\mathrm{feat}}.
\end{cases}
\]
The output is an encoded feature matrix used only for sequence modeling.

\subsection{Standardization for the CNN}

The sequence features are standardized with empirical mean and standard deviation computed from the designated sequence-training rows. If $z$ denotes a raw encoded feature dimension and $\mu,\sigma$ denote its training mean and standard deviation, then
\[
\tilde z = \frac{z - \mu}{\sigma}.
\]
If $\sigma$ is numerically degenerate, it is replaced by $1$.

\subsection{Sequence Construction}

For each company $j$ and each eligible month $t$, the sequence tensor is
\[
\seq_{j,t} =
\begin{bmatrix}
\tilde x_{j,t-L+1} \\
\tilde x_{j,t-L+2} \\
\vdots \\
\tilde x_{j,t}
\end{bmatrix}^{\top}
\in \R^{d \times L},
\]
where $L = \texttt{SEQ\_LEN}$. The transpose produces channel-first input for the CNN.

\subsection{Residual CNN Training}

The sequence model solves
\[
\hat g = \argmin_{g \in \mathcal{F}_{\mathrm{CNN}}} \sum_{(\seq_i,r_i) \in \mathcal{S}_{\mathrm{train}}} (r_i - g(\seq_i))^2.
\]
The final predictor is
\[
\hat y_i = \hat y^{\mathrm{CB}}_i + \hat g(\seq_i).
\]

\section{Training Objectives and Evaluation}

\subsection{CatBoost Objective}

CatBoost is fit with RMSE loss. The empirical objective is
\[
\loss_{\mathrm{CB}}(f) = \frac{1}{n} \sum_{i=1}^n (y_i - f(x_i))^2.
\]

\subsection{CNN Objective}

The CNN is fit with mean squared error on residuals:
\[
\loss_{\mathrm{CNN}}(g) = \frac{1}{m} \sum_{i=1}^m (r_i - g(\seq_i))^2.
\]

\subsection{Regression Metrics}

For any prediction vector $\hat y$ and truth vector $y$ over a finite evaluation set:
\[
\mathrm{RMSE}(y,\hat y) = \sqrt{\frac{1}{n} \sum_{i=1}^n (y_i - \hat y_i)^2},
\]
\[
\mathrm{MAE}(y,\hat y) = \frac{1}{n} \sum_{i=1}^n |y_i - \hat y_i|.
\]
The coefficient of determination is
\[
R^2 = 1 - \frac{\sum_{i=1}^n (y_i - \hat y_i)^2}{\sum_{i=1}^n (y_i - \bar y)^2},
\quad \bar y = \frac{1}{n}\sum_{i=1}^n y_i.
\]

\subsection{Directional Accuracy}

Directional accuracy is the proportion of sign agreement:
\[
\mathrm{DA}(y,\hat y) = \frac{1}{n} \sum_{i=1}^n \1\{\operatorname{sign}(y_i)=\operatorname{sign}(\hat y_i)\}.
\]

\subsection{Monthly Information Coefficient}

For each month $t$, the code computes the rank correlation between true and predicted returns. If $\mathrm{rank}(\cdot)$ denotes within-month ranking and $\rho$ denotes Pearson correlation, then
\[
\mathrm{IC}(t) = \rho\big(\mathrm{rank}(y_{\cdot,t}),\mathrm{rank}(\hat y_{\cdot,t})\big).
\]

\subsection{Top-K Return Metric}

For each month, the top $K$ predictions are selected and their realized returns averaged. If $\mathcal{I}_{t,K}$ denotes the index set of the $K$ largest predicted returns at time $t$, then
\[
\mathrm{TopK}(t) = \frac{1}{K} \sum_{i \in \mathcal{I}_{t,K}} y_{i,t}.
\]

\subsection{Residual Diagnostics}

Let $r^{\mathrm{base}}$ denote the CatBoost residual and $r^{\mathrm{final}}$ the final residual after CNN correction. Variance reduction is
\[
\mathrm{VR} = 1 - \frac{\Var(r^{\mathrm{final}})}{\Var(r^{\mathrm{base}})}.
\]
The one-step company-level autocorrelation is computed by grouping by company and using lag one within each group.

\section{Pseudocode}

The following pseudocode appears before the appendices and abstracts the operational logic of the notebook.

\begin{verbatim}
INPUT: train.csv, validation.csv, test.csv

1. Read each split and attach a split label.
2. Concatenate the three splits temporarily.
3. Remove stale engineered columns.
4. Convert Month to datetime, co_code to integer, target to numeric.
5. Normalize categorical labels as strings and fill missing categories.
6. Sort by (co_code, Month).
7. Collapse duplicate (co_code, Month) rows by keeping the last row.
8. For each company:
     a. compute lag returns for 1, 3, and 6 months,
     b. compute rolling mean and rolling standard deviation,
     c. compute lagged log-return values when available,
     d. compute month_index,
     e. add missingness indicators.
9. Split the processed table back into train, validation, and test.
10. Build CatBoost feature matrices and fit CatBoost on train.
11. Predict on validation and test; compute base residuals.
12. Partition validation into CNN-train and CNN-holdout months.
13. Encode sequence features and standardize them on training statistics.
14. Build causal windows of length SEQ_LEN for eligible rows.
15. Fit the CNN on residual sequences.
16. Predict residual corrections for validation and test sequences.
17. Form final predictions as CatBoost prediction + CNN correction.
18. Compute RMSE, MAE, R^2, IC, directional accuracy, and diagnostics.
\end{verbatim}

\appendix

\section{CatBoost from First Principles}

This appendix develops CatBoost-style gradient boosting from the beginning. The notation is intentionally explicit.

\subsection{Supervised Regression as Risk Minimization}

Let $\{(x_i,y_i)\}_{i=1}^n$ be a training set with $x_i \in \R^d$ and $y_i \in \R$. The regression problem is to find a function $f$ that minimizes the expected risk
\[
\mathcal{R}(f) = \E[\ell(Y,f(X))].
\]
Because the data distribution is unknown, we minimize the empirical risk
\[
\hat{\mathcal{R}}(f) = \frac{1}{n} \sum_{i=1}^n \ell(y_i,f(x_i)).
\]
For squared loss, $\ell(y,f(x))=(y-f(x))^2$.

\subsection{Function Space and Additive Modeling}

CatBoost builds an additive model
\[
F_M(x) = \sum_{m=0}^M \eta_m h_m(x),
\]
where each $h_m$ is a decision tree and $\eta_m$ is a step size. The model starts from a constant baseline
\[
F_0(x) = \argmin_c \sum_{i=1}^n (y_i-c)^2 = \bar y.
\]

\subsection{Gradient Boosting Step}

At boosting iteration $m$, the model has current prediction $F_{m-1}$. The negative gradient of the squared loss is
\[
g_i^{(m)} = -\left.\frac{\partial}{\partial F}\ell(y_i,F)\right|_{F=F_{m-1}(x_i)} = y_i - F_{m-1}(x_i).
\]
Thus the negative gradient is exactly the residual.

The next tree $h_m$ is chosen to approximate these residuals:
\[
h_m \approx \argmin_{h \in \mathcal{H}} \sum_{i=1}^n \big(g_i^{(m)} - h(x_i)\big)^2.
\]
Then
\[
F_m(x) = F_{m-1}(x) + \eta_m h_m(x).
\]

\subsection{Decision Trees as Partition Functions}

A regression tree partitions feature space into leaves $\{R_1,\dots,R_K\}$. On each leaf the tree is constant:
\[
h(x) = \sum_{k=1}^K c_k \1\{x \in R_k\}.
\]
The parameters are the regions $R_k$ and the leaf values $c_k$.

\subsection{Why Trees Work Well Here}

The panel contains nonlinear interactions, threshold effects, and missingness structure. Trees naturally represent these via axis-aligned splits. A split such as
\[
x_j \le \tau
\]
divides observations into two groups, allowing separate leaf predictions.

\subsection{Squared Loss and Leaf Values}

If the tree structure is fixed, the optimal leaf value in region $R_k$ under squared loss is the mean residual on that leaf:
\[
c_k^* = \frac{1}{|I_k|} \sum_{i \in I_k} g_i^{(m)},
\]
where $I_k = \{i : x_i \in R_k\}$.

\subsection{Categorical Variables}

CatBoost is designed to handle categorical features without requiring manual one-hot encoding. At an abstract level, a categorical variable $c \in \mathcal{C}$ is mapped into a representation usable by tree splits, often through target-statistic-based encodings or ordered statistics that reduce leakage. The essential mathematical point is that a categorical feature is not treated as an ordered real number unless the encoding makes that order explicit.

The notebook supplies categorical columns directly to CatBoost. Therefore the model internally performs the category handling required to optimize split quality while preserving leakage control.

\subsection{Why Ordered Statistics Matter}

If a category encoding uses target information, then naive averaging can leak the label into the feature. CatBoost addresses this by using permutations or ordered schemes so that the encoding for a row is computed from earlier rows only. This is consistent with the pipeline's leakage-free principle.

\subsection{Optimization Objective in the Notebook}

The notebook uses RMSE as the evaluation criterion and trains the model with the CatBoost regression objective. In formal terms, the fitted model approximates
\[
\hat f_{\mathrm{CB}} = \argmin_{f \in \mathcal{F}_{\mathrm{tree}}} \frac{1}{n} \sum_{i=1}^n (y_i - f(x_i))^2.
\]
The GPU/CPU choice changes implementation, not the mathematical objective.

\section{CNN from First Principles}

This appendix develops the convolutional residual learner from the beginning.

\subsection{Discrete Sequence Data}

The CNN consumes a tensor
\[
S_i \in \R^{d \times L},
\]
where $d$ is the number of channels and $L$ is the sequence length. Each column corresponds to one time step, and each row corresponds to one feature channel.

\subsection{One-Dimensional Convolution}

Let $W \in \R^{d \times k}$ be a filter of width $k$ and let $b \in \R$ be a bias term. For a valid temporal index $\tau$, the pre-activation at location $\tau$ is
\[
z_\tau = b + \sum_{c=1}^d \sum_{u=0}^{k-1} W_{c,u} S_{c,\tau+u}.
\]
This is the discrete convolution operator specialized to one-dimensional time.

\subsection{Multiple Filters and Feature Maps}

If there are $K$ filters, then the layer output is a family of feature maps
\[
z^{(1)},\dots,z^{(K)}.
\]
Each filter detects a different temporal pattern. One filter may respond to momentum-like patterns, another to reversals, another to volatility bursts.

\subsection{Nonlinearity}

After convolution, the model applies ReLU:
\[
\sigma(u) = \max(0,u).
\]
Thus the activated feature map is
\[
a_\tau = \sigma(z_\tau).
\]
The nonlinearity allows the CNN to represent non-affine functions of the input sequence.

\subsection{Pooling}

The notebook uses adaptive average pooling to reduce each channel to a single scalar. If $a^{(k)}$ is the activated output of filter $k$, then the pooled representation is
\[
p_k = \frac{1}{|\Omega_k|} \sum_{\tau \in \Omega_k} a^{(k)}_\tau,
\]
where $\Omega_k$ is the set of positions retained by the pooling operator.

\subsection{Linear Readout}

The pooled vector $p \in \R^K$ is passed through a linear map
\[
\hat r = w^\top p + b.
\]
This is the residual prediction.

\subsection{Why Convolution Fits the Problem}

The residual correction should depend on local time patterns rather than global rearrangements. Convolution is appropriate because it is translation-equivariant along time: the same temporal pattern can be recognized regardless of its location within the window.

\subsection{Backpropagation and Training}

The CNN parameters are trained by minimizing MSE on residuals:
\[
\hat\loss_{\mathrm{CNN}} = \frac{1}{m} \sum_{i=1}^m (r_i - \hat r_i)^2.
\]
Gradient descent updates parameters in the direction of decreasing loss. If $\theta$ denotes all CNN parameters, then the generic update is
\[
\theta \leftarrow \theta - \alpha \nabla_\theta \hat\loss_{\mathrm{CNN}},
\]
with learning rate $\alpha > 0$.

\subsection{Input Channels in the Notebook}

The notebook augments the sequence channels with the CatBoost prediction. Mathematically, if $x_{j,t}$ is the base engineered vector, then the CNN input at time $t$ includes
\[
\tilde x_{j,t} = \big[x_{j,t};\hat y^{\mathrm{CB}}_{j,t}\big].
\]
This allows the CNN to learn a direct correction to the first-stage forecast.

\section{Recursive Definitions}

We define the core objects recursively so that every higher-level object is built from lower-level objects.

\subsection{Recursive Definition of a Company History}

For a fixed company $j$, define the history up to time $t$ recursively by
\[
H_{j,1} = r_{j,1},
\]
and for $t \ge 2$,
\[
H_{j,t} = (H_{j,t-1}, r_{j,t}).
\]
This history object is the minimal causal input for feature construction.

\subsection{Recursive Definition of Lagged Features}

Let $L_\ell(y)$ denote the lag-$\ell$ operator. Then
\[
L_1(y)_{j,t} = y_{j,t-1},
\]
and recursively,
\[
L_{\ell+1}(y)_{j,t} = L_\ell(y)_{j,t-1}.
\]
This makes each lag a repeated application of the same primitive shift.

\subsection{Recursive Definition of the Final Prediction}

The final prediction is built from the base learner and the residual learner:
\[
\hat y^{(0)}_{j,t} = \hat y^{\mathrm{CB}}_{j,t},
\]
\[
\hat y^{(1)}_{j,t} = \hat y^{(0)}_{j,t} + \hat e^{\mathrm{CNN}}_{j,t}.
\]
This is the final output of the pipeline.

\section{Interpretation}

The notebook implements a structured decomposition of the target into a tabular component and a temporal correction. CatBoost addresses heterogeneous categorical and numerical signals in the cross-section. The CNN addresses serial dependence in the residuals. The preprocessing pipeline is the mechanism that makes this decomposition mathematically valid by preserving chronology, generating causal features, and enforcing unique company-month keys.

\section{Conclusion}

We have defined the problem, the axioms, the preprocessing operators, the model decomposition, and the training and evaluation criteria in a formal notation. The appendices derived the CatBoost and CNN components from first principles so that the entire notebook can be read as a sequence of mathematical transformations rather than as a collection of ad hoc code blocks.






\section{Notation Table}

\begin{longtable}{>{\raggedright\arraybackslash}p{0.10\textwidth}
                  >{\raggedright\arraybackslash}p{0.23\textwidth}
                  >{\raggedright\arraybackslash}p{0.14\textwidth}
                  >{\raggedright\arraybackslash}p{0.20\textwidth}
                  >{\raggedright\arraybackslash}p{0.28\textwidth}}
\toprule
Symbol & Meaning & Type / Set & Notes & Physical Interpretation \\
\midrule
\endfirsthead

\toprule
Symbol & Meaning & Type / Set & Notes & Physical Interpretation \\
\midrule
\endhead

$j$ & Company index & Integer & Identifies a firm in the panel & A specific company (e.g., \texttt{co\_code}) in the dataset \\
$t$ & Month index & Integer / time & Chronological index & A specific calendar month (e.g., "2020-01-01") for a firm \\
$J$ & Number of companies & Positive integer & Total cross-sectional size & Total number of unique companies/stocks in the dataset \\
$T_j$ & Number of observed months for company $j$ & Positive integer & Can vary across firms & Number of months with data available for this specific company \\
$y_{j,t}$ & Observed monthly gross return & Real number & Prediction target & The actual return achieved by the stock in month $t$ \\
$r_{j,t}$ & Raw observation record & Data object & Pre-engineering row-level data & A single row for company $j$ at time $t$ from the raw CSV \\
$x_{j,t}$ & Engineered tabular feature vector & $\mathbb{R}^d$ & CatBoost input & A row from the dataset containing all encoded features \\
$S_{j,t}$ & Sequence tensor & $\mathbb{R}^{d \times L}$ & CNN input built from past features & A sequence of the past $L$ months of feature rows for a company \\
$L$ & Sequence length & Positive integer & Window size used by CNN & The number of historical months in the CNN window \\
$d$ & Number of features / channels & Positive integer & Feature dimension after encoding & Total number of input variables/columns per row \\
$f$ & Full predictor & Function & Combined model & The entire ML pipeline predicting stock returns \\
$f_{\mathrm{CB}}$ & CatBoost predictor & Function & First-stage tabular learner & The initial tree model predicting from a flat feature vector \\
$f_{\mathrm{CNN}}$ & CNN residual predictor & Function & Second-stage temporal learner & The 1D CNN predicting target corrections \\
$\hat y_{j,t}$ & Final prediction & Real number & CatBoost plus CNN correction & The pipeline's ultimate guess of the firm's return \\
$\hat y^{\mathrm{CB}}_{j,t}$ & CatBoost prediction & Real number & First-stage output & CatBoost's initial guess of the return \\
$\hat e^{\mathrm{CNN}}_{j,t}$ & CNN residual prediction & Real number & Second-stage output & The CNN's predicted correction to CatBoost's error \\
$e_{j,t}$ & Residual & Real number & $y_{j,t} - \hat y^{\mathrm{CB}}_{j,t}$ & The actual error CatBoost made against the real return \\
$\D^{\mathrm{raw}}$ & Raw dataset & Set of records & Preprocessing input & The unprocessed CSV datasets (\texttt{train.csv}, etc.) \\
$\D$ & Processed dataset & Set of records & Feature-engineered panel & The fully engineered DataFrame containing all features \\
$\D_s$ & Split $s$ dataset & Set of records & $s \in \{\mathrm{train},\mathrm{validation},\mathrm{test}\}$ & A specific subset of the data like the training set \\
$\phi$ & Feature construction map & Function & Uses only current and past information & The code logic deriving features from raw data \\
$\psi$ & Sequence construction map & Function & Maps panel rows to sequence tensors & The logic stacking sequential rows into PyTorch tensors \\
$\mathcal{T}$ & Time domain & Ordered set & Calendar time support & The set of all calendar months represented in the panel \\
$\mathcal{F}_{\mathrm{CB}}$ & CatBoost hypothesis class & Function class & Set of tabular predictors & All possible tree ensembles CatBoost could learn \\
$\mathcal{F}_{\mathrm{CNN}}$ & CNN hypothesis class & Function class & Set of sequence predictors & All possible weight configurations for the CNN \\
$\loss_{\mathrm{CB}}$ & CatBoost loss & Scalar objective & Squared error objective & The Root Mean Squared Error (RMSE) minimized by CatBoost \\
$\loss_{\mathrm{CNN}}$ & CNN loss & Scalar objective & MSE on residuals & The Mean Squared Error (MSE) minimized by the CNN \\
$\hat{\mathcal{R}}(f)$ & Empirical risk & Scalar objective & Average training loss & The average error calculated over the actual training rows \\
$\R$ & Real numbers & Set & Continuous values & Continuous numerical values \\
$\N$ & Natural numbers & Set & Positive integers & Integer counts \\
$\1$ & Indicator function & Map & $\1\{\cdot\}$ equals 1 when condition holds & A binary/boolean flag (e.g., missing value indicators) \\
$\Var(\cdot)$ & Variance & Scalar statistic & Used in residual diagnostics & Statistical spread of the relative return errors \\
$\E[\cdot]$ & Expectation & Operator & Used in risk formulation & The expected value or population average \\
$\alpha$ & Learning rate & Positive real number & CNN optimization step size & Step size for the optimizer updating CNN weights \\
$w$ & Linear readout weights & Vector & Final CNN head parameters & The weights of the final linear layer in the CNN \\
$b$ & Linear readout bias & Scalar & Final CNN head intercept & The intercept of the final linear layer in the CNN \\
$W$ & Convolution kernel weights & Tensor & CNN filter parameters & Filter weights finding temporal patterns \\
$z_\tau$ & Convolution pre-activation & Real number & Output before nonlinearity & The raw CNN unit output before applying the ReLU activation \\
$a_\tau$ & Activation value & Real number & Output after ReLU & The CNN unit output after applying ReLU \\
$p_k$ & Pooled feature value & Real number & Adaptive average pooling output & The metric summarizing a CNN temporal feature map \\
$R_k$ & Tree leaf region & Subset of feature space & CatBoost/tree partition cell & A region defined by tree splits grouping similar firms \\
$c_k$ & Leaf value & Real number & Constant prediction on a leaf & The predicted return assigned to all firms in leaf $R_k$ \\
$g_i^{(m)}$ & Negative gradient / residual at boosting step $m$ & Real number & CatBoost boosting target & Unexplained return that the next tree attempts to predict \\
$F_m$ & Boosted ensemble after step $m$ & Function & Additive tree ensemble & The combined return prediction of the first $m$ trees \\
$H_{j,t}$ & Company history object & Sequence object & Recursive history representation & All known data for a company up to a given month \\
$\mathrm{lag\_ret}_\ell$ & Lagged return feature & Real number & $\ell$-month lag & The actual gross return of the stock $\ell$ months prior \\
$\mathrm{log\_ret\_lag}_\ell$ & Lagged log-return feature & Real number & Lagged log-return & The actual log-return of the stock $\ell$ months prior \\
$\mathrm{rolling\_mean}_w$ & Rolling mean feature & Real number & Window size $w$ & The average return of the stock over the last $w$ months \\
$\mathrm{rolling\_std}_w$ & Rolling standard deviation feature & Real number & Window size $w$ & Volatility of the stock's returns over the last $w$ months \\
$\mathrm{month\_index}$ & Calendar index & Integer & Year-month encoding & Integer indicating absolute calendar time (e.g., $2020 \times 12 + 1$) \\
$\mathrm{is\_missing}_f$ & Missingness indicator & $\{0,1\}$ & Records whether feature $f$ is missing & Flag indicating if a specific feature was originally missing \\
$\mathrm{RMSE}$ & Root mean squared error & Scalar metric & Regression evaluation & Typical error magnitude of the model's return predictions \\
$\mathrm{MAE}$ & Mean absolute error & Scalar metric & Regression evaluation & Average absolute magnitude of the return prediction errors \\
$R^2$ & Coefficient of determination & Scalar metric & Variance explained & Proportion of return variance explained by the model \\
$\mathrm{DA}$ & Directional accuracy & Scalar metric & Sign agreement rate & Percentage of months the predicted return sign was correct \\
$\mathrm{IC}(t)$ & Monthly information coefficient & Scalar metric & Rank correlation within month & Correlation between predicted and actual stock rankings \\
$\mathrm{TopK}(t)$ & Top-$K$ monthly realized return & Scalar metric & Portfolio-style metric & Average true return of stocks predicted to perform best \\
$\mathrm{VR}$ & Variance reduction & Scalar metric & Residual improvement diagnostic & The degree to which CNN reduced CatBoost's error variance \\
\bottomrule
\end{longtable}



\end{document}